% !PRINTED 751
\weberpage{751}

%<<CONT
schieden ist. Diese Zahl $\eta$, die doch auch dem Intervalle $(0, 1)$ angeh\"ort, kann mit keiner Zahl der Reihe $\Omega$ \"ubereinstimmen.

Man kann die Bildung von $\eta$ noch dadurch verallgemeinern, dass man die ersten $\beta$ bis zu einem beliebig weit entfernten willk\"urlich annimmt, und erst von da an das Gesetz: $\beta_\nu$ verschieden von $\alpha_\nu^{(\nu)}$, gelten l\"asst.

Da nun bewiesen ist, dass die reellen algebraischen Zahlen eine abz\"ahlbare Menge bilden, und dass es in jedem Intervalle Zahlen giebt, die einer gegebenen abz\"ahlbaren Menge nicht angeh\"oren, so folgt hieraus ganz unmittelbar:

\begin{center}
\gesperrt{Es giebt in jedem reellen Intervalle transcendente Zahlen.}
\end{center}

\begin{center}
\S.\,205.

\gesperrt{Transcendenz der Zahl $e$.}
\end{center}

Eine weit schwierigere Aufgabe ist es nun, von einer bestimmt vorgelegten Zahl zu entscheiden, ob sie algebraisch oder transcendent ist. Hier hat sich das Interesse haupts\"achlich auf die beiden in der Analysis so h\"aufig vorkommenden Zahlen $e$, d.\,h. die Basis des nat\"urlichen Logarithmensystems, und die \gesperrt{Ludolph}'sche Zahl $\pi$, das Verh\"altniss des Kreisumfanges zum Durchmesser, concentrirt.

F\"ur die Zahl $e$ ist die Frage von \gesperrt{Hermite} in einer ber\"uhmten Abhandlung entschieden, die f\"ur die sp\"ateren Untersuchungen \"uber die Zahl $\pi$ die Grundlage geworden ist$^{1)}$. Die Entscheidung f\"ur die Zahl $\pi$, die wegen ihrer Beziehung zu dem altber\"uhmten Problem der Quadratur des Kreises ganz besonders interessant war, bot aber noch lange Zeit un\"uberwindliche Schwierigkeiten. Endlich ist von \gesperrt{Lindemann} der Nachweis gef\"uhrt, dass auch $\pi$ zu den transcendenten Zahlen geh\"ort. Der von \gesperrt{Lindemann} gegebene Beweis bot aber dem Verst\"andniss zun\"achst noch grosse Schwierigkeiten, die durch sp\"atere Untersuchungen von \gesperrt{Weierstrass}, \gesperrt{Hilbert}, \gesperrt{Hurwitz} und \gesperrt{Gordan}$^{2)}$

\noindent\rule{4cm}{0.4pt}

$^{1)}$ \gesperrt{Hermite}, Sur la fonction exponentielle. Comptes rendus, T.\,LXXVII, 1873.

$^{2)}$ \gesperrt{Lindemann}, Ueber die Zahl $\pi$. Mathem. Annalen, Bd. 20, 1882. \gesperrt{Weierstrass}, Zu \gesperrt{Lindemann}'s Abhandlung \glqq Ueber die \gesperrt{Ludolph}'sche
%CONT>>
